%!LW recipe=pdflatex
\documentclass[a4paper,UKenglish,cleveref, autoref, thm-restate]{lipics-v2021}
%This is a template for producing LIPIcs articles. 
%See lipics-v2021-authors-guidelines.pdf for further information.
%for A4 paper format use option "a4paper", for US-letter use option "letterpaper"
%for british hyphenation rules use option "UKenglish", for american hyphenation rules use option "USenglish"
%for section-numbered lemmas etc., use "numberwithinsect"
%for enabling cleveref support, use "cleveref"
%for enabling autoref support, use "autoref"
%for anonymousing the authors (e.g. for double-blind review), add "anonymous"
%for enabling thm-restate support, use "thm-restate"
%for enabling a two-column layout for the author/affilation part (only applicable for > 6 authors), use "authorcolumns"
%for producing a PDF according the PDF/A standard, add "pdfa"

%\pdfoutput=1 %uncomment to ensure pdflatex processing (mandatatory e.g. to submit to arXiv)
%\hideLIPIcs  %uncomment to remove references to LIPIcs series (logo, DOI, ...), e.g. when preparing a pre-final version to be uploaded to arXiv or another public repository

%\graphicspath{{./graphics/}}%helpful if your graphic files are in another directory

\Crefname{section}{Sec.}{Secs.}
\Crefname{figure}{Fig.}{Figs.}

\bibliographystyle{plainurl}% the mandatory bibstyle

\usepackage[dvipsnames]{xcolor}
\usepackage{stmaryrd}
\usepackage{mathtools}
\usepackage{macros}
\setlength\fboxsep{2pt}

\title{Proof-Relevant Interpolation, Formalised}
  % TODO I'm sure we can do better

%\titlerunning{Dummy short title}

\author{Meven {Lennon-Bertrand}}{Université Paris Cité, INRIA, CNRS, IRIF}{meven.bertrand@inria.fr}{https://orcid.org/0000-0002-7079-8826}{}

\author{Alexis Saurin}{Université Paris Cité, INRIA, CNRS, IRIF}{Alexis.Saurin@irif.fr}{https://orcid.org/0009-0002-1304-5518}{}

\authorrunning{M. Lennon-Bertrand and A. Saurin}

\Copyright{Meven Lennon-Bertrand and Alexis Saurin}

\ccsdesc[100]{\textcolor{red}{Replace ccsdesc macro with valid one}} %TODO mandatory: Please choose ACM 2012 classifications from https://dl.acm.org/ccs/ccs_flat.cfm 

\keywords{Interpolation, Bidirectional Typing, Lambda Calculus}

% \category{} %optional, e.g. invited paper

% \relatedversion{} %optional, e.g. full version hosted on arXiv, HAL, or other respository/website
%\relatedversiondetails[linktext={opt. text shown instead of the URL}, cite=DBLP:books/mk/GrayR93]{Classification (e.g. Full Version, Extended Version, Previous Version}{URL to related version} %linktext and cite are optional

%\supplement{}%optional, e.g. related research data, source code, ... hosted on a repository like zenodo, figshare, GitHub, ...
%\supplementdetails[linktext={opt. text shown instead of the URL}, cite=DBLP:books/mk/GrayR93, subcategory={Description, Subcategory}, swhid={Software Heritage Identifier}]{General Classification (e.g. Software, Dataset, Model, ...)}{URL to related version} %linktext, cite, and subcategory are optional

%\funding{(Optional) general funding statement \dots}%optional, to capture a funding statement, which applies to all authors. Please enter author specific funding statements as fifth argument of the \author macro.

% \acknowledgements{I want to thank \dots}%optional

%\nolinenumbers %uncomment to disable line numbering


%Editor-only macros:: begin (do not touch as author)%%%%%%%%%%%%%%%%%%%%%%%%%%%%%%%%%%
\EventEditors{John Q. Open and Joan R. Access}
\EventNoEds{2}
\EventLongTitle{42nd Conference on Very Important Topics (CVIT 2016)}
\EventShortTitle{CVIT 2016}
\EventAcronym{CVIT}
\EventYear{2016}
\EventDate{December 24--27, 2016}
\EventLocation{Little Whinging, United Kingdom}
\EventLogo{}
\SeriesVolume{42}
\ArticleNo{23}
%%%%%%%%%%%%%%%%%%%%%%%%%%%%%%%%%%%%%%%%%%%%%%%%%%%%%%

\begin{document}

\maketitle

%TODO mandatory: add short abstract of the document
\begin{abstract}
Interpolation is a cool property in logic. It can be strengthened to a
proof-relevant version, which takes proof terms into account. We give a new
proof of this proof-relevant interpolation in the simply-typed lambda-calculus
with sums, building on ideas from the type theory community, and formalise
it in \Rocq.
\end{abstract}

\section{Introduction}

\subparagraph*{Interpolation for the \stlcp}
\begin{itemize}
  \item Interpolation is cool
  \item It can be made proof-relevant, cf \cite{Cubric1994} for \stlcp
\end{itemize}

\subparagraph*{Interpolation meets normal forms}

\begin{itemize}
  \item The original proof in Čubrić~\cite{Cubric1994} is… not very clean
  \item We can, however, discern that it follows the structure of
    normal forms
  \item By making this characterisation explicit and inductive, we obtain a much
    clearer proof
\end{itemize}

\subparagraph*{Normalisation for commutative cuts}

The key difficulty is that \(\beta\)-normal forms are not enough: interpolation
works well for terms/proofs which satisfy the subformula property, but in
\stlcp \(\beta\)-normal forms do not have this, we need to go further and
normalise with respect to \emph{commutative cuts}.

\subparagraph*{Contributions}
\begin{itemize}
  \item We develop the meta-theory of \stlcp, in particular normalisation
    for a calculus with commutative cuts
  \item We use this as a basis for a new proof of proof-relevant interpolation
    for the \stlcp
  \item All our work is formalised in \Rocq
\end{itemize}

\subparagraph*{A note on the meta-theory} We work in dependent type theory as
a meta-theory. To avoid confusion between object-level and meta-level
judgments, we denote meta-level typing as \(\mty\), and use \(\Set\) and \(\Prop\)
respectively for the meta-level universes of types and propositions –
we will not need more than one universe.
\mlb{Do we need to introduce more notations?}

\section{Syntax of the simply-typed \(\lambda\)-calculus with sums}

\paragraph*{Types and terms}
We start with our version of the simply-typed \(\lambda\)-calculus with sums,
hereafter abbreviated as \stlcp,
which we take to contain function types \(\to\), product types \(\times\),
sum types \(+\), and unit and empty types \(\top\) and \(\bot\).
To be able to state an interesting version of interpolation, we must also
include an unspecified set of base types.
To make the result more interesting,
we likewise include a number of constants and equations.
\mlb{Should this go here, or in a separate section?}

\begin{definition}[Language]
  \label{def:lang}
  A \emph{language} \(\lang\) consists of:
  \begin{itemize}
    \item a set of base types \(\base \in \Set\);
    \item a set \(\const \in \Set\) of constants;
    \item types for each constant, \ie
      a function \(\tyof \in \const \to \Ty_{\base}\).%
    \footnote{\(\Ty_{\base}\) is the set of types generated by
    \(\base\), to be defined in \cref{fig:stlc}.} 
  \end{itemize}
  We write \(\emptyset\) for the empty language, \ie that with no base types nor
  constants.
\end{definition}

\begin{definition}[Types and terms]
The definition of types and terms is given in \cref{fig:stlc}.
If \(\lang = (\base,\const,\tyof)\) we write
\(\Ty_{\lang}\) for \(\Ty_{\base}\), and similarly for contexts.
When it is not particularly relevant, we omit the language subscript.
\end{definition}

\begin{figure}
\begin{mathpar}
  \mprset{andskip=1em plus 0.5fil minus 0.5em}
  \vspace*{-1em}
  \jform{\ctxty[\base]{\Gamma}}[Contexts \(\Gamma,\Delta,\dots \in \Ctx_{\base}\)]
  %
  \inferdef{ }{\ctxty{\emp}} \and
  \inferdef{\ctxty{\Gamma} \\ \ctxty{A}}{\ctxty{\Gamma . (x : A)}} \\
  %
  \jform{\ctxty[\base]{A}}[Types \(A,B,\dots \in \Ty_{\base}\)]
  %
  \inferdef[Base]{b \in \base}{\ctxty[\base]{b}} \and
  \inferdef[Unit]{ }{\ctxty{\top}} \and
  \inferdef[Prod]{\ctxty{A} \\ \ctxty{B}}{\ctxty{A \times B}} \and
  \inferdef[Fun]{\ctxty{A} \\ \ctxty{B}}{\ctxty{A \to B}} \\
  \inferdef[Empty]{ }{\ctxty{\bot}} \and
  \inferdef[Sum]{\ctxty{A} \\ \ctxty{B}}{\ctxty{A + B}} \\
  %
  \jform{\typing[\lang]{\Gamma}{t}[A]}[Terms \(t,u,\dots \in \Tm_{\lang}(\Gamma,A)\)]
  \inferdef[Cst]{c \in \const}{\typing[(\base,\const,\tyof)]{\Gamma}{c}[\tyof(c)]} \and
  \inferdef[Var]{(x : A) \in \Gamma}{\typing{\Gamma}{x}[A]} \and
  \inferdef[Star]{ }{\typing{\Gamma}{\st}[\top]} \and
  \inferdef[Pair]{\typing{\Gamma}{t}[A] \\ \typing{\Gamma}{u}[B]}{
      \typing{\Gamma}{(t,u)}[A \times B]} \and
  \inferdef[Proj1]{\typing{\Gamma}{p}[A \times B]}{
      \typing{\Gamma}{\proj{1}{p}}[A]} \and
  \inferdef[Proj2]{\typing{\Gamma}{p}[A \times B]}{
      \typing{\Gamma}{\proj{2}{p}}[B]} \and
  \inferdef[Lam]{\typing{\Gamma.(x : A)}{t}[B]}{
      \typing{\Gamma}{\l x.t}[A \to B]} \and
  \inferdef[App]{\typing{\Gamma}{t}[A \to B] \\ \typing{\Gamma}{u}[A]}{
      \typing{\Gamma}{t~u}[B]} \and
  \inferdef[Raise]{\typing{\Gamma}{t}[\bot]}{\typing{\Gamma}{\rai t}[A]} \and
  \inferdef[InL]{\typing{\Gamma}{a}[A]}{\typing{\Gamma}{\inl a}[A + B]} \and
  \inferdef[InR]{\typing{\Gamma}{b}[B]}{\typing{\Gamma}{\inr b}[A + B]} \and
  \inferdef[If]{
    \typing{\Gamma}{s}[A + B] \\
    \typing{\Gamma.(x:A)}{b_l}[T] \\
    \typing{\Gamma.(x:B)}{b_r}[T] \\
  }{\typing{\Gamma}{\ite{s}{b_l}{b_r}}[T]}
%
\end{mathpar}
\caption{Types and terms of \stlcp with respect to a language
  \(\lang = (\base,\const,\tyof)\)}
\label{fig:stlc}
\end{figure}

\paragraph*{Substitutions}

We work with parallel substitutions, the action of which we denote
\(\subs{t}{\sigma}\).
Although we use variable names in terms for clarity,
we rather think them in the de Bruijn style \cite{deBruijn1972}, and
this is the way they are formalised in \Rocq. Accordingly, the
substitution calculus we use is essentially name-free, see \cref{fig:subst} for
the basic primitives. The typing judgment \(\typing{\Delta}{\sigma}[\Gamma]\)
means that \(\sigma\) take a variable \((x : A) \in \Gamma\) to a
term \(\sigma(x) \in \Tm(\Delta,A)\). The primitive's types should be already
rather informative, but let us describe them succinctly: \(\id\) is the identity
substitution, mapping each variable to itself; \(\sigma \pcomp \tau\) is composition,
mapping a variable \(x\) to \(\subs{\sigma(x)}{\tau}\); \(\wk\) drops the last
variable in context, and maps all other variables to themselves;%
\footnote{In a named syntax this is a no-op, but we include it
here as it changes a term is considered.} \(\sigma,t\) extends
a substitution \(\sigma\) by mapping the last variable in context to \(t\);
\(\Up \sigma\) maps the last variable to itself, and acts as \(\sigma\) otherwise;
finally \(t..\) is the substitution that maps the last variable to \(t\) and
acts as the identity otherwise. The following substitution rule is admissible:
\begin{mathpar}
  \inferdef[Subst]{\typing{\Delta}{\sigma}[\Gamma] \\ \typing{\Gamma}{\sigma}[A]}
    {\typing{\Delta}{\subs{t}{\sigma}}[A]}
\end{mathpar}

\begin{figure}
\begin{mathpar}
  \jform{\typing{\Delta}{\sigma}[\Gamma]}[Substitutions
    \(\sigma,\tau,\dots \in \Sub_{\lang}(\Delta,\Gamma)\)]
  \inferdef[Id]{ }{\typing{\Gamma}{\id}[\Gamma]} \and
  \inferdef[Id]{\typing{\Gamma}{\sigma}[\Delta] \and \typing{\Delta}{\tau}[\Theta]}{
      \typing{\Gamma}{\sigma \pcomp \tau}[\Theta]} \and
  \inferdef[Wk]{ }{\typing{\Gamma.(x:A)}{\wk}[\Gamma]} \and
  \inferdef[Cons]{\typing{\Delta}{\sigma}[\Gamma] \\ \typing{\Delta}{t}[A]}{
      \typing{\Delta}{\sigma,t}[\Gamma.(x:A)]} \and
  \inferdef[Up]{\typing{\Delta}{\sigma}[\Gamma]}{\typing{\Delta.(x:A)}{\Up \sigma}[\Gamma.(x:A)]} \and
  \inferdef[One]{\typing{\Gamma}{t}[A]}{\typing{\Gamma}{t..}[\Gamma.(x:A)]}
%
\end{mathpar}
\caption{Substitutions}
\label{fig:subst}
\end{figure}

\paragraph*{Reduction and equational theory}

The language defined in \cref{fig:stlc}, quotiented by the adequate equational
theory, is the internal language of bi-cartesian closed categories, \ie categories
with finite products, finite coproducts, and exponentials. However, studying
this equational theory is quite unwieldy due to the \(\eta\) laws for sums and
the empty type (see \eg Scherer~\cite{Scherer2017}). Thus, we concentrate
on a less powerful equational theory, which contains only the equations known
as “commutative cuts” \cite{??}. These are simpler to handle, and suffice to
obtain our ultimate result – see \cref{note:eta-sums,interpolation-eta}.
\mlb{For now, we don't allow extra equations, but we might.}

To factor out commuting conversions, we use a notion of eliminations.
\begin{definition}[Eliminations]
\emph{Eliminations} are given by the following grammar:
\[ e \Coloneqq \sapp{t} \mid \sproj{i} \mid \site{t}{t}\]
We define the \emph{plugging} of a term in an elimination \(\plug{e}{t}\) in the obvious way.
\end{definition}

Reductions, then, are of two kinds: the usual \(\beta\)-reductions, corresponding
to computation ---an abstraction on an argument, the projection of a pair,
and branching on a constructor--- and \emph{commuting conversions}, which “push”
an elimination across another one, the latter of a positive type (\ie \(\bot\) or
a sum). In the case of \(\bot\), the elimination gets completely removed, as it
is in a sense unnecessary: a proof of \(\bot\) can be directly used to inhabit
any type, there is no need to further eliminate it. In the case of the sum,
the elimination gets duplicated in each branch of the if.
\begin{definition}[Reductions]
  \emph{One-step reduction} \(\oredop\)
  is the congruent closure of the rules given in \cref{fig:stlc-red}.
  \emph{Reduction} \(\redop\) is the reflexive, transitive closure of one-step reduction.
\end{definition}

\begin{figure}
\begin{mathpar}
  % \jform{\conv{\Gamma}{t}{u}[A]}[Equality of terms]
  % \inferdef[Refl]{\typing{\Gamma}{t}[A]}{\conv{\Gamma}{t}{t}[A]} \and
  % \inferdef[Sym]{\conv{\Gamma}{t}{u}[A]}{\conv{\Gamma}{u}{t}[A]} \and
  % \inferdef[TRans]{\conv{\Gamma}{t}{u}[A] \\ \conv{\Gamma}{t}{v}[A]}
  %   {\conv{\Gamma}{t}{v}[A]} \\
  % \text{congruence/substitution rule} \\
  % \inferdef[Eq]{(t,u) \in \mathcal{E}(A)}{\conv{\Gamma}{t}{u}[A]} \and
  % \inferdef[Unit\(\eta\)]{\typing{\Gamma}{t}[\top]}{\conv{\Gamma}{t}{\st}[\top]} \\
  % \inferdef[Pair\(\beta\)]{\typing{\Gamma}{t}[A] \\ \typing{\Gamma}{u}[B]}{
  %   \conv{\Gamma}{\proj{1}{(t,u)}}{t}[A] \and
  %   \conv{\Gamma}{\proj{2}{(t,u)}}{u}[A]
  % } \and
  % \inferdef[Pair\(\eta\)]{\typing{\Gamma}{p}[A \times B]}{
  %   \conv{\Gamma}{p}{(\proj{1}{p},\proj{2}{p})}[A \times B]} \\
  % \inferdef[Fun\(\beta\)]{\typing{\Gamma,x:A}{t}[B] \\ \typing{\Gamma}{u}[A]}{
  %   \conv{\Gamma}{(\l x.t)~u}{\subs{t}{x}{u}}[A]} \and
  % \inferdef[Fun\(\eta\)]{\typing{\Gamma}{f}[A \to B]}{
  %   \conv{\Gamma}{f}{\l x.(f~x)}[A \to B]}
  \jform{\ored{t}{u}}[One-step reduction]
  \inferdef[Pair\textsubscript{1}\(\beta\)]{ }{\ored{\proj{1}{(t,u)}}{t}} \and
  \inferdef[Pair\textsubscript{2}\(\beta\)]{ }{\ored{\proj{2}{(t,u)}}{u}} \and
  \inferdef[Fun\(\beta\)]{ }{\ored{(\l x.t)~u}{\subs{t}{u..}}} \and
  \inferdef[InL\(\beta\)]{ }{\ored{\ite{\inl a}{b_l}{b_r}}{\subs{b_l}{a..}}} \and
  \inferdef[Inr\(\beta\)]{ }{\ored{\ite{\inr b}{b_l}{b_r}}{\subs{b_r}{b..}}} \\
  \inferdef[CommRais]{ }{\ored{\plug{e}{\rai t}}{\rai t}} \and
  \inferdef[CommIf]{ }{\ored{\plug{e}{\ite{s}{b_l}{b_r}}}{\ite{s}{\plug{e}{b_l}}{\plug{e}{b_r}}}}
\end{mathpar}
\caption{Reduction (congruence rules omitted)}
\label{fig:stlc-red}
\end{figure}

\mlb{We need to define the equational theory? But without confluence, there is
a priori a difference between the typed one and the untyped one, what do we do?}

\section{Bidirectional typing, normal forms and the subformula property}

In logic, a key characteristic of derivations is the \emph{subformula property}, satisfied
when where every formula appearing in the derivation is a subformula of some formula
in the conclusion judgment ---either as a hypothesis or conclusion. This
does not hold of arbitrary derivations, but is satisfied by derivations which
are sufficiently normalised, \ie cut-free proofs \cite{Girard1989}.
In a derivation satisfying the subformula property, no formula is ever “invented”,
a powerful structural property with many nice consequences. In
particular, interpolation is usually much easier to prove for cut-free derivation
\cite{??}. One typically obtains interesting properties for all provable
statements by a combination of a cut elimination theorem, which says that every
proof is equal to a cut-free one, and an analysis of cut-free proofs, which can
exploit its more constrained structure, including the subformula property.

For a type theorist, this might seem a bit mysterious. Cut elimination, by
the Curry-Howard correspondence, is well-known to coincide with reduction of
terms viewed as proof witnesses. However, what does the subformula property
amount to? This is much less known, but it can be very helpfully linked to
\emph{bidirectional typing} \cite[Sec. 8]{Dunfield2021}.

In bidirectional typing,
which originates from type-checking algorithms \cite{Coquand1996,Pierce2000},
the typing judgment is decomposed into the two, mutually defined
\emph{inference} and \emph{checking} judgments. The operational intuition is
that in the former the type is an output,%
\footnote{\textit{I.e.} the question is ``What is the type of this term?''.}
while in the latter it is an input.%
\footnote{And the question is ``Does this term have this type?''.}
This attention to the flow of information is, in a sense, also what
the subformula property intends to capture: in both cases, the core property
is that type information always comes from somewhere and is never invented.
And in fact bidirectional typing can helpfully be leveraged to exploit
this subformula property, see \eg the strengthening proof in
Lennon-Bertrand~\cite{LennonBertrand2021}.

With this view it is perhaps unsurprising that bidirectional
type systems can be used to structurally characterise normal forms, since,
as we just explained, these correspond to cut-free proofs which have a good
subformula property. This is, however, a remark with powerful consequences:
because bidirectional type systems have a nice, inductive structure, they
give a good basis on which to build clean proof of theorems that rely on
the subformula property, such as interpolation. This is the core idea which
we will follow in the remainder of this article. The rest of this section is
dedicated to setting up the bidirectional characterisation of normal forms and

\subsection{A bidirectional characterisation of normal forms}

\mlb{
Where we give the bidirectional type system, and prove it captures well-typed
normal forms.
}

\subsection{Normalisation for \stlcp}

\mlb{
Where we give a proof by logical relations that every term reduces to one in normal form.
}

\section{Interpolation}

\mlb{
Where we explain the core of our interpolation proof.
}

\section{Comments on the formalisation}

\mlb{Where we explain things on the \Rocq proof}
\begin{enumerate}
  \item Autosubst/Sulfur
  \item How far can we push automation? We do auto-loading + setoid rewriting, can
    we go further?
  \item Confluence proof?
\end{enumerate}

%%
%% Bibliography
%%

%% Please use bibtex, 

\bibliography{biblio}

\appendix

\end{document}
